\section{怎样发现并修复错误证明}
\label{sec:01-Mathematical-Language/03-Proof-Techniques/04}

先读一段推导。设 \(a=b\ne0\)：

\[
\begin{aligned}
a^2 &= ab,\\
a^2-b^2 &= ab-b^2,\\
(a-b)(a+b) &= b(a-b),\\
a+b &= b,\\
2b &= b,\\
2 &= 1.
\end{aligned}
\]

六行，每一行单看都像正当的恒等变形，结论却是 \(2=1\)。

如果你扫了两遍还是觉得处处合理，那正说明了这类错误的危险之处：它不靠语法错误暴露自己，而是藏在某一步操作的合法前提里。式子写得越流畅，越没人去问那一步凭什么能做。

\subsection{出错的不是代数，是被违反的前提}

问题出在第三行到第四行：从 \((a-b)(a+b)=b(a-b)\) 消掉 \(a-b\)。而题目开头就设了 \(a=b\)，所以 \(a-b=0\)。

除以零不是``不推荐''，是没有定义。一旦某一步不合法，从那里往后整条推导链就断了，后面写得再工整也不产生任何效力；末尾的 \(2=1\) 不是数学出了矛盾，而是这步非法操作的报警声。

由此得到审阅证明的第一原则：不要只问``这一行像不像代数变形''，要问``这一步用的是哪条规则，它的前提满足了吗''。消去律的前提是被消掉的因子非零，开方的前提是被开方数非负，取对数的前提是真数为正，交换求和与积分次序的前提是某种收敛性。规则都写在教科书里，被跳过的从来是前提。

同一原则的另一副面孔是不可逆操作。解 \(\sqrt{x+3}=x-3\) 时两边平方，得到 \(x=1\) 或 \(x=6\)。平方这步本身完全合法，但它只保证单向蕴含，不保证等价——原方程的解一定在这两个数里，这两个数却未必是原方程的解。回代一下：\(x=6\) 时两边都等于 3，通过；\(x=1\) 时左边为 2 而右边为 \(-2\)，出局。凡是用过平方、取绝对值、乘以一个可能为零的量之类的动作，候选解都必须回代验证，否则增根就混进了答案。

\subsection{以特例充普遍}

验证了 \(n=1,2,\dots,100\) 全都成立，能不能宣布对所有自然数成立？不能。

例子的用途是发现规律、建立猜想、找出反例，它提供信心而不提供证明。``验过一百个''和``验过全部''之间隔着一道全称量词，这道缝不会因为把一百换成一百万而合上。一段程序跑过一百万组输入没崩，也不能说明它对所有合法输入安全——边界值、整数溢出、浮点舍入的罕见组合，未必落在那一百万组里。

这是数据工作里最常见的一种不可靠结论。在一个数据集上验证过的规律，成立范围就是那个数据集的分布；把它说成``这个方法有效''，等于把一次特例检验当成了全称证明。要跨过这道缝，要么用归纳、要么用针对输入域结构的论证，要么老老实实把结论的适用范围写进句子里。

\subsection{箭头掉头与自我循环}

由 \(P\to Q\) 和 \(Q\) 成立，能不能推出 \(P\)？不能，这是肯定后件。\(x=2\) 蕴含 \(x^2=4\)，但知道 \(x^2=4\) 并不能断定 \(x=2\)，因为 \(-2\) 也在候选里。

自己写证明时特别容易在这儿栽：心里想要 \(P\)，于是顺着 \(P\) 往下推，推到一个明显成立的 \(Q\)，然后宣布 \(P\) 得证。可你证的是 \(P\to Q\)。检查办法很笨也很有效——给每一步标上箭头方向，再看整条链的方向是不是你需要的那个。

一个真实出现在作业里的例子：``若 \(n^2\) 是偶数，则 \(n\) 是偶数，因为偶数的平方是偶数。''后半句说的是 \(n\) 偶 \(\to\) \(n^2\) 偶，正好是要证方向的反面。修复的办法不是补字，是改走等价的逆否方向：设 \(n\) 为奇数，\(n=2k+1\)，则 \(n^2=2(2k^2+2k)+1\) 为奇数，于是原命题成立。方向歪了的时候，字数是帮不上忙的。

循环论证是同一种病的另一种表现：要证的结论被当成理由用掉了。``这两条直线平行，因为斜率相等；斜率相等，因为它们平行''——绕了一圈，始终没碰到定义、公理或任何已确认的事实。它的包装有时相当精巧，结论换个等价说法就认不出来了，比如要证 \(n\) 是偶数，中途用了一句``因为 \(n\) 能被 2 整除''。判别方法是把每条理由往回追，看它最终落在什么上面；如果追着追着回到了起点，那就是一个闭环。

\subsection{量词顺序被悄悄调换}

这一类最难自己发现，因为交换前后的两句话读起来几乎一样。

\begin{center}
\begin{tikzpicture}[x=1cm,y=1cm,font=\small,>=stealth]
  \node[black!70] at (1.3,3.4) {\(\forall x\,\exists y\ P(x,y)\)};
  \fill (0,2.4) circle (2pt); \node[anchor=east] at (-0.15,2.4) {\(x_1\)};
  \fill (0,1.4) circle (2pt); \node[anchor=east] at (-0.15,1.4) {\(x_2\)};
  \fill (0,0.4) circle (2pt); \node[anchor=east] at (-0.15,0.4) {\(x_3\)};
  \fill (2.6,2.4) circle (2pt); \node[anchor=west] at (2.75,2.4) {\(y_1\)};
  \fill (2.6,1.4) circle (2pt); \node[anchor=west] at (2.75,1.4) {\(y_2\)};
  \fill (2.6,0.4) circle (2pt); \node[anchor=west] at (2.75,0.4) {\(y_3\)};
  \draw[->,black!55] (0.15,2.4) -- (2.45,2.4);
  \draw[->,black!55] (0.15,1.4) -- (2.45,1.4);
  \draw[->,black!55] (0.15,0.4) -- (2.45,0.4);
  \node[black!70] at (1.3,-0.5) {各配各的};

  \node[black!70] at (7.9,3.4) {\(\exists y\,\forall x\ P(x,y)\)};
  \fill (6.6,2.4) circle (2pt); \node[anchor=east] at (6.45,2.4) {\(x_1\)};
  \fill (6.6,1.4) circle (2pt); \node[anchor=east] at (6.45,1.4) {\(x_2\)};
  \fill (6.6,0.4) circle (2pt); \node[anchor=east] at (6.45,0.4) {\(x_3\)};
  \fill (9.2,1.4) circle (2.6pt); \node[anchor=west] at (9.35,1.4) {\(y\)};
  \draw[->,black!55] (6.75,2.4) -- (9.05,1.5);
  \draw[->,black!55] (6.75,1.4) -- (9.05,1.4);
  \draw[->,black!55] (6.75,0.4) -- (9.05,1.3);
  \node[black!70] at (7.9,-0.5) {全体共用一个};
\end{tikzpicture}

\emph{两个量词只换了顺序，右边却强得多：它额外要求那个 \(y\) 不随 \(x\) 变。证出了左边而当作右边用，是错误证明里最难自查的一种。}
\end{center}

差别全在依赖关系上。左边允许 \(y\) 是 \(x\) 的函数，右边不允许——右边成立时左边自动成立，反过来不行。判断的窍门是问一句：你手里那个 \(y\) 是在见到 \(x\) 之前就固定下来的，还是见到 \(x\) 之后才造出来的？如果构造过程用到了 \(x\)，那你证到的只是左边。

调参的场景把这件事说得很直白。``对每个数据集，都存在一组超参数使模型达标''是一句相当弱的话，它允许每个数据集配一套自己的参数；``存在一组超参数，对所有数据集都达标''则强得多，那才是可以打包发布的东西。两句话中文只差了语序，工程价值差着量级。分析学里的逐点收敛与一致收敛、连续与一致连续，全是同一个陷阱的不同实例，后面读到那些定义时可以回来对照本图。相关的量词否定规则见 \hyperref[sec:01-Mathematical-Language/01-Logic/03]{``量词与否定''}。

\subsection{论域和边界被默认掉}

``若 \(ab=0\)，则 \(a=0\) 或 \(b=0\)。''这句话在实数和整数里成立，换到矩阵上就不成立——两个非零矩阵完全可以乘出零矩阵。结论没变，论域变了，真假就变了。

``连续函数一定取到最大值''同样漏了前提：定义域必须是有界闭区间。\(f(x)=x\) 在 \((0,1)\) 上连续，最大值和最小值都取不到。这类漏看有个规律——结论越眼熟，越容易漏掉它依赖的条件，因为大脑会自动把熟悉的句子补全成它最常出现的那个版本。

对付这类问题的办法是把定理当合同读：先找条件，再找结论，然后逐条核对眼前的对象是否落在条件里。条件不是装饰，每一条都在挡住某一类反例；说不出某个条件挡的是什么，通常意味着你还没真正读懂这条定理。

\subsection{换一双眼睛读自己的证明}

刚写完的时候是最不适合审阅的时刻——脑子里还留着大量没落到纸上的理由，看什么都觉得理所当然。隔几个小时再读效果会好很多；来不及等，就强行剥掉作者身份，假装这是别人交上来要你挑错的东西。

逐行拷问这几个问题：

\begin{enumerate}
\tightlist
\item
  对象和假设交代清楚了吗？变量属于哪个集合，量词的顺序是哪一种？
\item
  每一步蕴含的方向都是你需要的方向吗？有没有把只能正着用的定理反着用了？
\item
  有除法、开方、取对数或消去因子吗？它们各自的前提当场满足吗？
\item
  用过不可逆变换吗？候选解回代验证了吗？
\item
  分类讨论穷尽了吗？``当且仅当''的两个方向都走了吗？
\item
  说存在——见证给出来了，或者说明了凭什么必定存在吗？说唯一——两个候选相等的论证在哪里？
\item
  最后一句和原结论逐字比对过吗？有没有悄悄弱化成一句看起来很像的话？
\end{enumerate}

这份清单能拦下相当一部分错误，但要诚实说明它拦不住什么：有些错误来自对定理前提的深层误解，对着清单画勾是发现不了的，那时需要的是回去重读定理本身，而不是反复盯着自己的推导。还有一些错误只在特定的论域里才现形，清单里的每一条都答``是''，命题依然可以是假的。

修复错误证明的顺序也因此是固定的：先定位是哪一步不合法或哪一句被默认了，再判断这个漏洞是能补的还是致命的，最后才谈重写。跳过定位直接补字，通常只是把漏洞挪了个位置。

本章到此把``从已知走到结论''的几种基本走法讲完了。但有一类论域用穷尽切分也覆盖不了——全体自然数没法分成有限块。\hyperref[ch:008]{``归纳、递归与不变量：用局部规则控制无限过程''}处理的正是这种情况：用一条局部规则，换取对无限多个对象的控制。
